\documentclass[11pt,a4paper]{article}
\usepackage[T1]{fontenc}
\usepackage[utf8]{inputenc}
\usepackage{amsmath,amssymb,xcolor}
\usepackage[margin=25mm]{geometry}
\usepackage[colorlinks=true,urlcolor=blue]{hyperref}
\definecolor{nexorablue}{HTML}{2459B5}
\setlength{\parindent}{0pt}
\setlength{\parskip}{7pt}
\setlength{\emergencystretch}{2em}
\renewcommand{\familydefault}{\sfdefault}
\begin{document}

{\small\bfseries\color{nexorablue} NEXORA CPGE}\par
{\small MPSI / NEX-RGT-04 / Énoncé}\par
{\LARGE\bfseries Une bobine en deux phases}\par
\vspace{0.25em}\hrule\vspace{0.6em}

À courant initial nul, une bobine de $0{,}12\,\mathrm H$ est alimentée sous $12\,\mathrm V$ en série avec $6\,\mathrm\Omega$. Après $40\,\mathrm{ms}$, on retire la source et on ferme la boucle sur $9\,\mathrm\Omega$. Déterminer le courant à la commutation, puis la durée nécessaire pour tomber à $0{,}25\,\mathrm A$. Calculer l’énergie dissipée pendant toute la seconde phase.


\end{document}
